A scaffold can hide a racial price penalty from the three explainers most used in geographic machine learning. How
much it can hide is set by each explainer's estimator: LIME, whose queries are unlike real homes, is fooled almost for
free, while KernelSHAP and GeoShapley, bound by efficiency and querying many near-real hybrids, can be fooled only at a
cost in the bias actually delivered. Making an explainer spatially explicit does not change this. Under attack
GeoShapley's location effect shrinks and spurious interactions appear, so its decomposition should not be read as
evidence that a model treats places fairly. Audits of spatial models should ask where an explainer queries the model,
and should pair explanations with checks on outcomes that a model cannot route around.
